\begin{exercisechunk}
\item \textbf{Teaching, active learning, and passive learning.}
\label[exercise]{ex:l6-threshold-information}
\exerciseprofile{2}{2}{2}{\exproof\quad\excalculation\quad\exinterpretation}{%
Distinguish three information structures by deriving their requirements for
identifying a threshold.}
Let \(N\geq6\), let \(\mathcal K=\{2,\ldots,N-1\}\), and, for
\(k\in\mathcal K\), define \(f_k(x)=\I\{x\geq k\}\) on \([N]\).
The goal is to identify \(k\) exactly from labeled examples.
\begin{enumerate}[(a),beginpenalty=10000]
\item A teacher knows \(k\) and may choose labeled examples. Prove that
\(\{(k-1,0),(k,1)\}\) identifies \(k\) uniquely within \(\mathcal K\).
Explain why no single labeled example identifies every \(k\in\mathcal K\).
\item An active learner does not know \(k\), but may choose each query after
observing the previous labels. Give a binary-search strategy that identifies
every \(k\in\mathcal K\) using at most
\(\lceil\log_2(N-2)\rceil\) queries. Prove that no active strategy can have a
smaller worst-case number of binary-label queries.
\item A passive learner receives \(X_1,\ldots,X_n\) independently and uniformly
from \([N]\), together with their labels. Fix \(k\in\{3,\ldots,N-2\}\).
Prove that the observed labels identify \(k\) uniquely if and only if both
\(k-1\) and \(k\) occur in the sample, and that the probability of this event is
\[
 1-2\left(1-\frac1N\right)^n
   +\left(1-\frac2N\right)^n.
\]
Deduce that \(n\geq N\log(2/\delta)\) suffices for this probability to be at
least \(1-\delta\).
\item Explain why the two-example teaching set is not evidence that dependence
alone makes learning faster. State precisely what target information the
teacher has that the active and passive learners do not.
\end{enumerate}
\begin{hints}
For (b), maintain the set of thresholds consistent with all answers. For the
lower bound, a depth-\(q\) binary decision tree has at most \(2^q\) leaves. For
(c), use inclusion--exclusion and then a union bound.
\end{hints}
\end{exercisechunk}
